\section{MDL 把拟合与描述长度放在同一尺度}

一条高阶曲线可以穿过每个训练点，但``把所有点逐个记住''并没有发现可复用规律。Minimum description length 把这一直觉变成编码问题：若模型真的抓住结构，那么把模型告诉接收者之后，剩余数据应更容易编码。

最简单的 two-part MDL 选择最小化

\[
L(M)+L(D\mid M)
\]

的模型。第一项编码模型结构和参数，第二项在已知模型后编码数据或残差。复杂模型通常减小第二项，却增大第一项；最好的解释是在总码长上取得平衡。

\subsection{Likelihood 为什么能变成码长}

概率模型给事件 \(D\) 的理想码长

\[
L(D\mid M)\approx-\log_2p(D\mid M).
\]

底数取 2 即以 bit 为单位（第一节已说明 log 底只改变数值单位）。``理想''与 \(\approx\) 是因为真实码长必须为整数：逐符号编码与 \(-\log_2p\) 相差不到 1 bit，对长序列用算术编码可在平均意义下逼近它。对数据分布取期望，这正是第一节``用模型分布编码真实来源''的平均码长，即 cross-entropy。

高 likelihood 意味着模型把短码分配给实际出现的数据。但一个模型若能通过大量自由参数把任何数据都赋高 likelihood，接收者还需要知道这些参数，不能只算残差码长。

prefix code 的长度满足 Kraft inequality：

\[
\sum_x2^{-L(x)}\le1.
\]

这一不等式来自 prefix code 对应的树结构，两个方向的严格讨论见``前缀码与Kraft不等式''。它保证码长可以对应某个 sub-probability；反过来，概率分布也能诱导接近 \(-\log_2p(x)\) 的可译码。MDL 因而不是给损失随意套上``bit''这个单位，而是在概率、编码和模型比较之间建立操作性联系。

\subsection{复杂度不等于参数个数}

连续参数需要指定精度或使用 universal code。两个参数数量相同的模型，如果参数可表达范围、可识别性和对数据的灵敏度不同，编码成本也可能不同。神经网络中大量参数对称还会让``逐个参数编码''重复计算等价函数。

BIC 在正则、固定维度、大样本等条件下可由边缘似然或某些编码近似导出，但它不是所有 MDL 方法的同义词。奇异模型、混合模型和深度网络可能违反其 Hessian 非退化等假设。更精细的 normalized maximum likelihood、Bayesian mixture code 和 prequential code 对复杂度有不同定义与有限样本行为。

prequential 编码特别直观：先用前一段数据训练模型，再用模型给下一段数据编码，持续更新。一个迅速学会规律的模型会较早给未来数据短码；一个只能靠记忆的模型在早期要付出高预测代价。此时训练算法本身也成为编码方案的一部分。

\subsection{一个传感器曲线的例子}

考虑用多项式解释温度与传感器漂移。阶数上升时，训练残差负 log-likelihood 持续下降；但还要编码阶数、系数及其精度。低阶模型可能残差较长，高阶模型参数码过长，总码长通常在中间出现最低点。

这一选择仍依赖噪声模型和编码方案。若残差重尾却使用 Gaussian code，少数异常会被赋予极长码；换成稳健噪声模型可能比盲目提高多项式阶数更短。MDL 迫使我们把``模型结构''和``解释不了的部分''都说清，却不会自动替我们选择唯一合理的 model class。

实际比较时应固定数据精度、log 底和可解码协议，分别报告结构、参数与数据码长。用 held-out 或 prequential code 做敏感性检查，并确认改变参数量化精度后结论是否稳定。

\subsubsection{把拟合曲线拆成两张账单}

为一组带噪曲线依次拟合 0 到 12 阶多项式。画出训练负 log-likelihood、一个明确 two-part 参数码长和总码长。改变参数量化精度与噪声分布，观察最优阶数如何变化。最后说明为什么``训练 NLL 最小''\,``BIC 最小''和``你定义的总码长最小''不必给出同一个模型。
